\chapter{Langage, fonctions et calcul : les bases}
\section*{Exercice 1.1 -- Nier une affirmation}
La négation est : il existe un réel $x$ tel que, pour tout entier $n$, $n\le x$. L'affirmation initiale est la propriété d'Archimède, vraie dans $\R$ (chapitre 36). Changer l'ordre des quantificateurs changerait le sens.

\section*{Exercice 1.2 -- Domaine avant calcul}
Le logarithme impose $x>1$ et la racine au dénominateur impose $x<4$. Domaine : $]1,4[$. Sur cet intervalle, le dénominateur est non nul, donc $f(x)=0\Leftrightarrow\ln(x-1)=0\Leftrightarrow x=2$.

\section*{Exercice 1.3 -- Une équivalence à corriger}
Il faut $x\ge0$. Le carré donne $x^2-x-6=(x-3)(x+2)=0$. La racine $-2$ ne vérifie pas le signe initial ; seul $3$ convient. Le carré conserve une égalité, mais sa réciproque demande deux membres de même signe.

\section*{Exercice 1.4 -- Composer}
$(g\circ f)(x)=x^2+1$ et $(f\circ g)(x)=(x+1)^2=x^2+2x+1$. Elles coïncident seulement en $x=0$, et ne sont donc pas égales comme fonctions. La composition n'est pas commutative.

\section*{Exercice 1.5 -- Bijection et restriction}
Sur $[0,+\infty[$ vers $[0,+\infty[$, chaque $y\ge0$ a un unique antécédent positif ou nul $\sqrt y$. La réciproque est donc la racine carrée. Sur $\R$, les nombres opposés non nuls ont la même image, empêchant l'injectivité.

\section*{Exercice 1.6 -- Égalité d’ensembles}
Pour $x\in E$, appartenir au membre de gauche signifie $f(x)\in A\cup B$, donc $f(x)\in A$ ou $f(x)\in B$. Cela équivaut à l'appartenance au membre de droite. L'équivalence pour tout $x$ prouve les deux inclusions.
